% !TeX root = ../report_project_2stage_OTA.tex

\section{Conclusion}

This project developed a low-power two-stage CMOS OTA using the SKY130A
open-source PDK. The design progressed from first-order calculations and
an ideal two-stage model to a transistor-level implementation. This
approach made it possible to study the gain, pole--zero behavior, and
Miller compensation independently before introducing device limitations
and parasitic effects.

The final OTA operates from a nominal $1.8~\mathrm{V}$ supply and draws
approximately $4.0~\mu\mathrm{A}$, including the reference bias branch.
With the physical compensation network
$C_{\mathrm{MILLER}}=736.44~\mathrm{fF}$ and
$R_{\mathrm{MILLER}}=180~\mathrm{k}\Omega$, the nominal open-loop DC gain is
$89.73~\mathrm{dB}$. At the specified $2~\mathrm{pF}$ nominal load, the
unity-gain frequency reaches $1.184~\mathrm{MHz}$ with a phase margin of
$89.04^\circ$. At the maximum $4~\mathrm{pF}$ load, the phase margin
remains $74.98^\circ$. All 45 simulated PVT conditions satisfy the hard
gain and stability requirements, with minimum values of $87.63~\mathrm{dB}$
DC gain, $64.69^\circ$ phase margin, and $13.76~\mathrm{dB}$ gain margin.

Statistical characterization also informed the design choices. Increasing
the area of the first-stage NMOS mirror loads reduced the nominal
input-referred offset standard deviation from $6.125~\mathrm{mV}$ to
$1.109~\mathrm{mV}$. Subsequent iterations showed diminishing returns,
while reducing the capacitive loading introduced by the second-stage NMOS
helped recover stability. The final nominal offset spread is
$3\sigma(V_{OS})=3.33~\mathrm{mV}$; the largest value among the nine
voltage--temperature conditions is $3.83~\mathrm{mV}$. Thus, the
approximately $3~\mathrm{mV}$ statistical-offset objective was not fully
met, despite the substantial improvement. The integrated output noise
was also reduced to $120.27~\mu\mathrm{V}_{\mathrm{RMS}}$ over
$1~\mathrm{Hz}$--$1~\mathrm{MHz}$.

The remaining limitations are primarily related to physical
implementation and verification. The reported results are schematic-level
simulations using the available SKY130A models; no layout, extracted
parasitics, or silicon measurements are included. A natural continuation
would be a matching-aware layout followed by DRC, LVS, parasitic
extraction, and post-layout characterization. This would allow the
stability, offset, noise, and usable input/output voltage ranges to be
reassessed with physical parasitics and layout-dependent effects.
